\subsection{Product of two measures}

Let S and T be two topological spaces, equipped respectively with two (positive) premeasures \(\lambda\) and \(\mu\) , and let X be the product space \(S \times T\) . Let K be a compact subset of X; let us denote by A and B the projections of K on S and T respectively, and set

\[
\nu_ {\mathrm{K}} = \left(\lambda_ {\mathrm{A}} \otimes \mu_ {\mathrm{B}}\right) _ {\mathrm{K}}.\tag{3}
\]

We thus define a premeasure on X. For, let L be a compact subset of X containing K, and let C and D be its two projections; then \(A \subset C\) , \(B \subset D\) , consequently, using the transitivity of induced measures, and Prop. 12 of Ch. V, §8, No.~5, we have

\[
\begin{array}{r l} & {(\nu_ {\mathrm{L}}) _ {\mathrm{K}} = \big ((\lambda_ {\mathrm{C}} \otimes \mu_ {\mathrm{D}}) _ {\mathrm{L}} \big) _ {\mathrm{K}} = (\lambda_ {\mathrm{C}} \otimes \mu_ {\mathrm{D}}) _ {\mathrm{K}}} \\ & {\qquad = \big ((\lambda_ {\mathrm{C}} \otimes \mu_ {\mathrm{D}}) _ {\mathrm{A} \times \mathrm{B}} \big) _ {\mathrm{K}} = (\lambda_ {\mathrm{A}} \otimes \mu_ {\mathrm{B}}) _ {\mathrm{K}} = \nu_ {\mathrm{K}}.} \end{array}
\]

DEFINITION 6. --- The premeasure \(\nu\) defined by (3) is called the product premeasure of \(\lambda\) and \(\mu\) , and is denoted \(\lambda \otimes \mu\) .

This definition obviously extends to the case that \(\lambda\) and \(\mu\) are complex premeasures, and one then has \(|\lambda\otimes\mu|=|\lambda|\otimes|\mu|\) (Ch. III, §4, No.~2, Prop. 3 and Ch. IV, §5, No.~7, Lemma 3).

We conserve the notations of Ch. III, §4 and Ch. V, §8 relative to products of measures and to iterated integrals. In particular, if \(f\) and \(g\) are two functions defined respectively on S and T, with values in \(\overline{\mathbf{R}}_+\) or in \(\mathbf{C}\) , the function \((s,t)\mapsto f(s)g(t)\) on \(S\times T\) will be denoted \(f\otimes g\) .

PROPOSITION 10. --- Let \(\nu\) be the product premeasure of \(\lambda\) and \(\mu\) ; for every function \(f \in \mathcal{F}_{+}(\mathrm{S})\) and every function \(g \in \mathcal{F}_{+}(\mathrm{T})\) ,

\[
\nu^ {\bullet} (f \otimes g) = \lambda^ {\bullet} (f) \mu^ {\bullet} (g).\tag{4}
\]

The premeasure \(\nu\) is the only premeasure on \(S \times T\) that satisfies (4).

As K (resp. L) runs over the set of compact subsets of S (resp. of T), we have

\[
\begin{array}{l} \nu^ {\bullet} (f \otimes g) = \sup _ {K, L} \nu_ {K \times L} ^ {\bullet} \big ((f \otimes g) _ {K \times L} \big) = \sup _ {K, L} (\lambda_ {K} \otimes \mu_ {L}) ^ {\bullet} (f _ {K} \otimes g _ {L}) \\ \qquad = \sup _ {K, L} \lambda_ {K} ^ {\bullet} (f _ {K}) \cdot \mu_ {L} ^ {\bullet} (g _ {L}) = \big (\sup _ {K} \lambda_ {K} ^ {\bullet} (f _ {K}) \big) \big (\sup _ {L} \mu_ {L} ^ {\bullet} (g _ {L}) \big) \\ \qquad = \lambda^ {\bullet} (f) \mu^ {\bullet} (g) \end{array}
\]

by Prop. 8 of Ch. V, §8, No.~3.

Let \(\eta\) be a second premeasure on \(S \times T\) satisfying (4), and let \(K\) and \(L\) be compact subsets of \(S\) and \(T\) respectively, \(f\) and \(g\) elements of \(\mathcal{F}_{+}(K)\) and \(\mathcal{F}_{+}(L)\) respectively. One has the relation \((f \otimes g)^{0} = f^{0} \otimes g^{0}\) between the extensions by 0, therefore (§1, No.~2, Prop. 2)

\[
\begin{array}{r l} & {\eta_ {\mathrm{K} \times \mathrm{L}} ^ {\bullet} (f \otimes g) = \eta^ {\bullet} \big ((f \otimes g) ^ {0} \big) = \eta^ {\bullet} (f ^ {0} \otimes g ^ {0})} \\ & {\qquad = \lambda^ {\bullet} (f ^ {0}) \mu^ {\bullet} (g ^ {0}) = \lambda_ {\mathrm{K}} ^ {\bullet} (f) \mu_ {\mathrm{L}} ^ {\bullet} (g).} \end{array}
\]

In particular, if one takes \(f \in \mathcal{K}_{+}(\mathrm{K})\) , \(g \in \mathcal{K}_{+}(\mathrm{L})\) , one sees that \(\eta_{\mathrm{K} \times \mathrm{L}}\) has the characteristic property of the product measure \(\lambda_{\mathrm{K}} \otimes \mu_{\mathrm{L}}\) (Ch. III, §4, No.~1, Th. 1). Therefore \(\eta_{\mathrm{K} \times \mathrm{L}} = \nu_{\mathrm{K} \times \mathrm{L}}\) ; since every compact subset of \(S \times T\) is contained in a set of the form \(K \times L\) , the transitivity of induced measures implies that \(\eta = \nu\) .

COROLLARY 1. --- If \(\lambda\) and \(\mu\) are measures, then \(\nu\) is a measure.

For, let \(x = (s, t)\) be a point of X, and let U and V be neighborhoods of s, t respectively, such that \(\lambda^{\bullet}(\mathrm{U}) < +\infty\) , \(\mu^{\bullet}(\mathrm{V}) < +\infty\) ; the set \(U \times V\) is a neighborhood of x, and \(\nu^{\bullet}(\mathrm{U} \times \mathrm{V}) = \lambda^{\bullet}(\mathrm{U}) \mu^{\bullet}(\mathrm{V}) < +\infty\) by (4); the encumbrance \(\nu^{\bullet}\) is thus locally bounded, and the premeasure \(\nu\) is a measure.

This result extends at once to complex measures.

COROLLARY 2. --- If A is a subset of S locally negligible for \(\lambda\) , then \(A \times T\) is locally \(\nu\) -negligible.

COROLLARY 3. --- Suppose that \(\lambda\) (resp. \(\mu\) ) is the sum of a summable family \((\lambda_{\alpha})_{\alpha \in A}\) (resp. \((\mu_{\beta})_{\beta \in B}\) ) of measures on S (resp. T). The family \((\lambda_{\alpha} \otimes \mu_{\beta})_{(\alpha, \beta) \in A \times B}\) is then summable, and its sum is \(\lambda \otimes \mu\) .

For, let \(p\) be the encumbrance \(\sum_{\alpha, \beta} (\lambda_{\alpha} \otimes \mu_{\beta})^{\bullet}\) ; if \(f \in \mathcal{F}_{+}(S)\) and \(g \in \mathcal{F}_{+}(T)\) , then obviously \(p(f \otimes g) = \lambda^{\bullet}(f)\mu^{\bullet}(g)\) . The proof of Cor. 1 then shows that \(p\) is locally bounded, so that the family \((\lambda_{\alpha} \otimes \mu_{\beta})\) is summable (§1, No.~7, Prop. 7). Its sum \(\eta\) then satisfies \(\eta^{\bullet} = p\) (§1, No.~7, Prop. 7), and Prop. 10 implies \(\eta = \nu\) .

